\scanpage{14}

\deleted{\uline{Proposition} Soit $A$ anneau de Dedekind, $G$ gr. \inserted{plat} affine \inserted{de t.f.} sur $A$, d'algèbre $B$ satisfaisant aux conditions du lemme 1. \deleted{Alors} Soit $N$ comme dans le lemme 2, et considérons \deleted{\unsure{la}} la représentation de $G$ dans $N$ induite par la repr. régulière gauche. Alors l'hom. $G\to\underline{\operatorname{Aut}}(N)$ est un morphisme fidèlement plat}

\figure{\begin{tikzpicture}[x=1cm,y=1cm]\draw (0,0)--(12,.75);\end{tikzpicture}}

\[
\left\{\begin{array}{ll}
\underline B=\varinjlim B_i &\\
1)\quad B/B_i\text{ plat} & \deleted{\text{donc}}\quad B_i\otimes B_i\subset B\otimes B\\
2)\quad \Delta B_i\subset B_i\otimes B_i &\\
3)\quad B_i\text{ loc. libre de t.f.} &
\end{array}\right.
\]
\deleted{\unsure{Suppos.} $X$ \uline{quasi-cpt} \uline{qu. sép.}}

Soit $E$ un $\underline O_X G$-module quasi-coh., \deleted{\illegible{1}}

Alors 1) Pour tt sous-module quasi-coh. $F$ il existe un plus petit \unsure{$A$-module|$\underline O_X$-module} quasi-coh. $F'$ de $F$ \deleted{tel que} stable par $G$, et sa formation commute à tt changement de base plat. Si $F$ est de t.f. $F'$ itou.

2) Si $X$ quasi-compact quasi-séparé, alors $E$ est $\varinjlim$ de sous-modules \inserted{qu. coh.} de t.f. qui sont stables par $G$.

\inserted{\deleted{\unsure{Pour $X$ quasi-cpt quasi-sép.}} \unsure{Si} $X$ quasi-cpt quasi-sép.}

3) Soit $F\to E$ un hom., avec $F$ quasi-coh. \inserted{de t.f.} dont l'image engendre $E$. Alors il existe

\mspage{13}
